\documentclass[10pt,a4paper]{article}

%double si on veut une version prof et une version élève. simple si on veut une seule version.
\newcommand{\buildmode}{simple}

\InputIfFileExists{version.tex}{}{}
% Version (définie par le script)
\providecommand{\version}{eleve}

\usepackage{feuilleinterro}

%%%à modifier à chaque fois

\renewcommand{\niveau}{2nde}
\renewcommand{\chapitre}{Proportions, pourcentages}
\renewcommand{\numchap}{0.2}
\renewcommand{\theme}{Statistiques}

%%%titre "CORRECTION" au lieu de "INTERROGATION", sans les champs Nom/Prénom
\renewcommand{\interroversion}[1]{
    \noindent\textbf{\niveau}

    \vspace{-0.5cm}
    \begin{center}
        \textbf{\large CORRECTION -- \chapitre}
    \end{center}

    \vspace{0.1cm}
}

\begin{document}

% =========================================================
% PAGE 1 : VERSION A
% =========================================================

\interroversion{A}

\textbf{Exercice 1 -- Tableau de valeurs \hfill /2}

On considère la fonction \(f\) définie sur \(\mathbb{R}\) par \(f(x)=x^2-2x-3\).

\[
\begin{aligned}
f(-1) &= (-1)^2-2\times(-1)-3 = 1+2-3 = 0\\
f(2)  &= 2^2-2\times 2-3 = 4-4-3 = -3
\end{aligned}
\]

\begin{center}
    \renewcommand{\arraystretch}{1.5}
    \setlength{\tabcolsep}{0.25cm}
    \begin{tabular}{|*{7}{c|}}
    \hline
    \(x\) & \(-2\) & \(-1\) & \(0\) & \(1\) & \(2\) & \(3\) \\
    \hline
    \(f(x)\) & \(5\) & \textcolor{blue}{\(0\)} & \(-3\) & \(-4\) & \textcolor{blue}{\(-3\)} & \(0\) \\
    \hline
    \end{tabular}
\end{center}

\textbf{Exercice 2 -- Pourcentages \hfill /3}

\begin{tasks}(2)
    \task Calculer \(10\%\) de \(90\).

    \(10\%\) de \(90\), c'est \(\dfrac{10}{100}\times 90\). \textcolor{blue}{
    \[
    \frac{10}{100}\times 90 = \frac{10\times 90}{100} = \frac{900}{100} = 9.
    \]}
    \task Poser et effectuer le calcul de \(35\%\) de \(80\).

    \textcolor{blue}{
    \[
    35\%\text{ de }80 = \frac{35}{100}\times 80 = \frac{35\times 80}{100} = \frac{2\,800}{100} = 28.
    \]
    (On pose \(35\times80\) : \(35\times80=35\times8\times10=280\times10=2\,800\), puis on divise par \(100\).)}
    \task Calculer \(20\%\) de \(15\%\) de \(300\).

    \textcolor{blue}{
    \[
    \begin{aligned}
    15\%\text{ de }300 &=\frac{15}{100}\times300=\frac{4\,500}{100}=45,\\
    20\%\text{ de }45 &=\frac{20}{100}\times45=\frac{900}{100}=9.
    \end{aligned}
    \]}
    \task Calculer le tiers de \(21\).

    \textcolor{blue}{Le tiers de \(21\), c'est \(21\div3=7\) (on vérifie : \(7\times3=21\)).}
\end{tasks}

\newpage

% =========================================================
% PAGE 2 : VERSION B
% =========================================================

\interroversion{B}

\textbf{Exercice 1 -- Tableau de valeurs \hfill /2}

On considère la fonction \(f\) définie sur \(\mathbb{R}\) par \(f(x)=x^2+x-2\).

\[
\begin{aligned}
f(-1) &= (-1)^2+(-1)-2 = 1-1-2 = -2\\
f(3)  &= 3^2+3-2 = 9+3-2 = 10
\end{aligned}
\]

\begin{center}
    \renewcommand{\arraystretch}{1.5}
    \setlength{\tabcolsep}{0.25cm}
    \begin{tabular}{|*{7}{c|}}
    \hline
    \(x\) & \(-2\) & \(-1\) & \(0\) & \(1\) & \(2\) & \(3\) \\
    \hline
    \(f(x)\) & \(0\) & \textcolor{blue}{\(-2\)} & \(-2\) & \(0\) & \(4\) & \textcolor{blue}{\(10\)} \\
    \hline
    \end{tabular}
\end{center}

\textbf{Exercice 2 -- Pourcentages \hfill /3}

\begin{tasks}(2)
    \task Calculer \(20\%\) de \(50\).

    \textcolor{blue}{
    \[
    \frac{20}{100}\times 50 = \frac{20\times 50}{100} = \frac{1\,000}{100} = 10.
    \]}
    \task Poser et effectuer le calcul de \(24\%\) de \(75\).

    \textcolor{blue}{
    \[
    24\%\text{ de }75 = \frac{24}{100}\times 75 = \frac{24\times 75}{100} = \frac{1\,800}{100} = 18.
    \]
    (On pose \(24\times75\) : \(24\times75=24\times70+24\times5=1\,680+120=1\,800\), puis on divise par \(100\).)}
    \task Calculer \(10\%\) de \(40\%\) de \(500\).

    \textcolor{blue}{
    \[
    \begin{aligned}
    40\%\text{ de }500 &=\frac{40}{100}\times500=\frac{20\,000}{100}=200,\\
    10\%\text{ de }200 &=\frac{10}{100}\times200=\frac{2\,000}{100}=20.
    \end{aligned}
    \]}
    \task Calculer le quart de \(28\).

    \textcolor{blue}{Le quart de \(28\), c'est \(28\div4=7\) (on vérifie : \(7\times4=28\)).}
\end{tasks}

\newpage

% =========================================================
% PAGE 3 : VERSION C
% =========================================================

\interroversion{C}

\textbf{Exercice 1 -- Tableau de valeurs \hfill /2}

On considère la fonction \(f\) définie sur \(\mathbb{R}\) par \(f(x)=x^2-3x+2\).

\[
\begin{aligned}
f(-2) &= (-2)^2-3\times(-2)+2 = 4+6+2 = 12\\
f(3)  &= 3^2-3\times 3+2 = 9-9+2 = 2
\end{aligned}
\]

\begin{center}
    \renewcommand{\arraystretch}{1.5}
    \setlength{\tabcolsep}{0.25cm}
    \begin{tabular}{|*{7}{c|}}
    \hline
    \(x\) & \(-2\) & \(-1\) & \(0\) & \(1\) & \(2\) & \(3\) \\
    \hline
    \(f(x)\) & \textcolor{blue}{\(12\)} & \(6\) & \(2\) & \(0\) & \(0\) & \textcolor{blue}{\(2\)} \\
    \hline
    \end{tabular}
\end{center}

\textbf{Exercice 2 -- Pourcentages \hfill /3}

\begin{tasks}(2)
    \task Calculer \(25\%\) de \(40\).

    \textcolor{blue}{
    \[
    \frac{25}{100}\times 40 = \frac{25\times 40}{100} = \frac{1\,000}{100} = 10.
    \]}
    \task Poser et effectuer le calcul de \(18\%\) de \(50\).

    \textcolor{blue}{
    \[
    18\%\text{ de }50 = \frac{18}{100}\times 50 = \frac{18\times 50}{100} = \frac{900}{100} = 9.
    \]
    (On pose \(18\times50\) : \(18\times50=18\times5\times10=90\times10=900\), puis on divise par \(100\).)}
    \task Calculer \(50\%\) de \(20\%\) de \(200\).

    \textcolor{blue}{
    \[
    \begin{aligned}
    20\%\text{ de }200 &=\frac{20}{100}\times200=\frac{4\,000}{100}=40,\\
    50\%\text{ de }40 &=\frac{50}{100}\times40=\frac{2\,000}{100}=20.
    \end{aligned}
    \]}
    \task Calculer le tiers de \(18\).

    \textcolor{blue}{Le tiers de \(18\), c'est \(18\div3=6\) (on vérifie : \(6\times3=18\)).}
\end{tasks}

\newpage

% =========================================================
% PAGE 4 : VERSION D
% =========================================================

\interroversion{D}

\textbf{Exercice 1 -- Tableau de valeurs \hfill /2}

On considère la fonction \(f\) définie sur \(\mathbb{R}\) par \(f(x)=x^2+2x-1\).

\[
\begin{aligned}
f(-1) &= (-1)^2+2\times(-1)-1 = 1-2-1 = -2\\
f(2)  &= 2^2+2\times 2-1 = 4+4-1 = 7
\end{aligned}
\]

\begin{center}
    \renewcommand{\arraystretch}{1.5}
    \setlength{\tabcolsep}{0.25cm}
    \begin{tabular}{|*{7}{c|}}
    \hline
    \(x\) & \(-2\) & \(-1\) & \(0\) & \(1\) & \(2\) & \(3\) \\
    \hline
    \(f(x)\) & \(-1\) & \textcolor{blue}{\(-2\)} & \(-1\) & \(2\) & \textcolor{blue}{\(7\)} & \(14\) \\
    \hline
    \end{tabular}
\end{center}

\textbf{Exercice 2 -- Pourcentages \hfill /3}

\begin{tasks}(2)
    \task Calculer \(10\%\) de \(70\).

    \textcolor{blue}{
    \[
    \frac{10}{100}\times 70 = \frac{10\times 70}{100} = \frac{700}{100} = 7.
    \]}
    \task Poser et effectuer le calcul de \(32\%\) de \(75\).

    \textcolor{blue}{
    \[
    32\%\text{ de }75 = \frac{32}{100}\times 75 = \frac{32\times 75}{100} = \frac{2\,400}{100} = 24.
    \]
    (On pose \(32\times75\) : \(32\times75=32\times70+32\times5=2\,240+160=2\,400\), puis on divise par \(100\).)}
    \task Calculer \(25\%\) de \(40\%\) de \(240\).

    \textcolor{blue}{
    \[
    \begin{aligned}
    40\%\text{ de }240 &=\frac{40}{100}\times240=\frac{9\,600}{100}=96,\\
    25\%\text{ de }96 &=\frac{25}{100}\times96=\frac{2\,400}{100}=24.
    \end{aligned}
    \]}
    \task Calculer le quart de \(36\).

    \textcolor{blue}{Le quart de \(36\), c'est \(36\div4=9\) (on vérifie : \(9\times4=36\)).}
\end{tasks}

\end{document}
